\question[4]
In einer Schulbibliothek werden Bücher verwaltet. Die Datenbank enthält die folgende Tabelle:

\begin{center}
\begin{tabular}{|c|l|l|c|}
\hline
BuchID & Titel & Autor & Erscheinungsjahr \\
\hline
101 & Die Physiker & Dürrenmatt & 1962 \\
102 & Tschick & Herrndorf & 2010 \\
103 & Krabat & Preußler & 1971 \\
\hline
\end{tabular}
\end{center}

Erläutere die Begriffe \textbf{Entität}, \textbf{Entitätstyp},
\textbf{Attribut}, \textbf{Tabelle}, \textbf{Datensatz} und
\textbf{Datenfeld} anhand des Beispiels.

\begin{solutionbox}{5cm}
\begin{itemize}
\item \textbf{Entität:} Ein konkretes Objekt der realen Welt, z.\,B. das Buch „Tschick“.
\item \textbf{Entitätstyp:} Die Klasse ähnlicher Objekte, hier „Buch“.
\item \textbf{Attribut:} Eine Eigenschaft eines Entitätstyps, z.\,B. Titel oder Autor.
\item \textbf{Tabelle:} Die strukturierte Darstellung eines Entitätstyps in Zeilen und Spalten.
\item \textbf{Datensatz:} Eine Zeile der Tabelle, die alle Informationen zu einer Entität enthält.
\item \textbf{Datenfeld:} Ein einzelner Wert innerhalb eines Datensatzes, z.\,B. „Tschick“ oder „2010“.
\end{itemize}
\end{solutionbox}